\section{The two-point skeleton and its moment defect}
\label{ms:geometry-counts}

Every centered law can be separated into its positive and nonpositive
parts. Replacing each part by its conditional mean gives a two-point law;
the remaining conditional fluctuations measure its departure from that
law. We first make this decomposition exact, then identify a moment
defect that vanishes at Esseen's extremal distribution. All quantities
in the construction belong to the same marginal $P$.

\subsection{The sign split and its units}

Fix $P\in\mathcal P$ and put
\[
 p=\Pp(X>0),\quad q=1-p,\quad v=pq,\quad
 m=\E X_+=\E X_->0,\quad B_0=\mathbf1_{\{X>0\}}.
\]
Both $p,q$ are positive; a possible atom at zero is assigned to the
nonpositive class. Define $V=\E[X\mid B_0]$ and $W=X-V$, so that
\[
 V=-m/q\quad(B_0=0),\qquad V=m/p\quad(B_0=1),\qquad
 \theta=\E W^2=1-m^2/v<1.
\]
Write $\theta_-=\E[W^2;B_0=0]$ and $\theta_+=\E[W^2;B_0=1]$.
The physical distance between the two conditional means is
\[
 \ell=m/v=\sqrt{(1-\theta)/v}>0.
\]
Set $U=W/\ell$ and $Y=B_0+U$. Then
\begin{equation}
 X=\ell(Y-p),\qquad \E[U\mid B_0]=0.
\label{ms:bernoulli-coordinates}
\end{equation}
Let $U_0,U_1$ have the two conditional laws, and put
\[
 w_i=\E U_i^2,\quad \mu_i=\E U_i^3,\quad
 s=qw_0+pw_1=\frac{v\theta}{1-\theta},\quad
 \lambda=ns,\quad \ell=(v+s)^{-1/2}.
\]
The moments $\mu_i$ are finite under the original third-moment assumption.
Higher moments used later require bounded support or a separate proved
moment bound. At a selected extremizer, bounded support is available from
Lemma~\ref{ms:contact}; a quantitative support radius remains a further step.

If $s=0$, both nonnegative conditional variances vanish, so $U_0=U_1=0$
almost surely. This is exactly the standardized Bernoulli law. If $s>0$,
the parameter $\alpha=pw_1/s$ lies in $[0,1]$ and gives
$w_0=s(1-\alpha)/q$, $w_1=s\alpha/p$. No division by $s$ is used on the
zero-noise face.

The signs of $X$ are still those defining the two classes, so expansion
of the actual conditional moments gives
\begin{align}
 \rho/\ell^3&=v(1-2v)+3v(w_0+w_1)-q\mu_0+p\mu_1,\notag\\
 \gamma/\ell^3&=v(q-p)+3v(w_1-w_0)+q\mu_0+p\mu_1,\notag\\
 \beta/\ell^2&=v(q-p)+pw_1-qw_0.
\label{ms:conditional-moments}
\end{align}
For example, on the nonpositive class $|X|^3=\ell^3(p-U_0)^3$;
on the positive class it is $\ell^3(q+U_1)^3$. Conditional centering removes
the linear terms and gives the first identity; the other two follow in
the same way. The common pair $(\mu_0,\mu_1)$ must be retained whenever
several such identities are used together.

\subsection{A nonnegative physical defect}

We now suppose the support is compact. Define
\[
 \tau(t)=\E[X\mathbf1_{\{X\ge t\}}],\qquad \sigma=\sqrt{10}.
\]
For $t<0$, centering gives $\tau(t)=-\E[X\mathbf1_{\{X<t\}}]\ge0$;
for $t\ge0$ nonnegativity is immediate. Signwise Tonelli integration gives
\begin{equation}
 \int_\R\tau(t)\,dt=1,
 \qquad 2\int_\R|t|\tau(t)\,dt=\rho.
\label{ms:tau-moments}
\end{equation}
On every bounded component $(a,b)$ of the complement of the closed support,
$\tau$ is constant in the interior; denote that value by $\tau_{a,b}$.
Set
\begin{equation}
 J=\sum_{(a,b)}\tau_{a,b}(b-a)^2,
 \qquad D=(3+\sigma)\rho-\gamma-3J.
\label{ms:physical-defect-definition}
\end{equation}
There are at most countably many gaps. If the support lies in $[-R,R]$,
then $J\le2R\int\tau=2R$, so this sum converges. Endpoint atoms do not
alter the Lebesgue integrals defining its weights.

\begin{lemma}[Physical defect and the central gap]
\label{ms:physical-defect}
\label{ms:fourier:afcore:lem:physical-defect}
The quantity $D$ in \eqref{ms:physical-defect-definition} is nonnegative.
If $0$ belongs to the support, then $D\ge(\sigma-1)\rho$. Consequently
$D<(\sigma-1)\rho$ implies a genuine central support gap $(-A,B)$ with
$A,B>0$.
\end{lemma}
\begin{proof}
If zero belongs to the support, every gap has one sign and
$(b-a)^2\le2\int_a^b|t|\,dt$. Sum using the constant gap weights and
\eqref{ms:tau-moments} to obtain $J\le\rho$. Thus
$D\ge\sigma\rho-\gamma\ge(\sigma-1)\rho$.

Otherwise compactness and the presence of both signs give the central
gap $(-A,B)$. Write
\[
 T_+=\E X_+^3-mB^2,\quad T_-=\E X_-^3-mA^2,
 \quad R_+=\E X_+^2-mB,\quad R_-=\E X_-^2-mA,
\]
all nonnegative. Let $J_o=J-m(A+B)^2$ be the sum over the outer gaps.
For a positive outer gap $(a,b)$,
\[
 (b-a)^2\le2\int_a^b(t-B)\,dt.
\]
Summing and using Tonelli bounds its total contribution by
\[
 J_o^+\le2\int_B^\infty(t-B)\tau(t)\,dt
       =\E[X_+(X_+-B)^2]=T_+-2BR_+\le T_+.
\]
Reflection of this integral calculation, without changing the fixed law,
gives $J_o^-\le T_--2AR_-\le T_-$. Expansion of $D$ yields the identity
\begin{align}
 D={}&m\{\sqrt{\sigma+1}A-\sqrt{\sigma-1}B\}^2
       +(\sigma-1)T_++(\sigma+1)T_-\notag\\
    &+3(T_++T_--J_o).\label{ms:defect-decomposition}
\end{align}
Here $\sqrt{\sigma+1}\sqrt{\sigma-1}=3$. Every term is nonnegative,
which proves the claim.
\end{proof}

\begin{corollary}[The zero-defect law]\label{ms:zero-defect}
For a compactly supported $P\in\mathcal P$, $D=0$ if and only if
\[
 X=(B_0-p_*)/\sqrt{p_*q_*},\qquad
 p_*=(4-\sqrt{10})/2,\quad q_*=1-p_*,
 \quad B_0\sim\operatorname{Bernoulli}(p_*).
\]
The sign of the third moment is fixed by the definition of $D$.
\end{corollary}
\begin{proof}
If $D=0$, zero cannot lie in the support, so there is a central gap
$(-A,B)$. The nonnegative terms in \eqref{ms:defect-decomposition}
force $T_+=T_-=0$. Since
$T_+=\E[X_+(X_+^2-B^2)]$ and the positive support lies in $[B,\infty)$,
all positive mass is at $B$; similarly all negative mass is at $-A$.
Centering and unit variance give $pB=qA$ and $AB=1$.
The square in \eqref{ms:defect-decomposition} now gives
\[
 \frac pq=\frac AB=\sqrt{\frac{\sigma-1}{\sigma+1}}
                   =\frac{\sigma-1}{3},
\]
hence $p=p_*$. Conversely this two-point law has $T_+=T_-=J_o=0$
and the square vanishes, so $D=0$.
\end{proof}

The defect therefore records the loss from the extremal two-point
geometry. For quantitative estimates we use the full nonnegative
decomposition, rather than only its equality case.

Several elementary consequences will be useful.
 Integrating $(X+A)(X-B)\ge0$ gives $AB\le1$.
Cauchy--Schwarz gives $1/(2\rho)\le m\le1/2$. The outer-gap estimates give
$J\le\rho+2m\le\rho+1$. Since $\sigma-1>2$,
\eqref{ms:defect-decomposition} implies $T_++T_-\le D/2$; if $D<2$,
then
\[
 m(A^2+B^2)=\rho-T_+-T_-\ge1-D/2,
 \qquad A+B\ge\sqrt{2-D}.
\]
These are necessary conditions of the same physical law, rather than
additional definitions of a relaxed defect.

\subsection{One joint budget for endpoint motion and noise}

Let $\bar A=m/q$, $\bar B=m/p$ be the conditional mean magnitudes, and
let $\Delta_A=\bar A-A$, $\Delta_B=\bar B-B$. They are nonnegative.
Define $p_*=(4-\sqrt{10})/2$.

\begin{lemma}[Simultaneous gap inequality]
\label{ms:simultaneous-gap}
For this central gap,
\begin{align}
 D\ge{}&\frac{(6+2\sigma)(1-\theta)^{3/2}}{\sqrt{pq}}(p-p_*)^2
 +6m\{(A+B)(\Delta_A+\Delta_B)+\Delta_A\Delta_B\}\notag\\
 &+3(\sigma+1)B\theta_++3(\sigma+3)A\theta_-.
\label{ms:simultaneous-gap-formula}
\end{align}
\end{lemma}
\begin{proof}
The stronger outer-gap bound in the preceding proof gives
\[
 D\ge(\sigma-1)\E X_+^3+(\sigma+1)\E X_-^3
                      -6mAB+6BR_++6AR_-.
\]
Conditional cubic expansion gives
\[
 \E X_+^3=m\bar B^2+
       \E[(X+2\bar B)(X-\bar B)^2;X>0]
       \ge m\bar B^2+3B\theta_+,
\]
and similarly $\E X_-^3\ge m\bar A^2+3A\theta_-$.
Moreover $R_+=\theta_++m\Delta_B$ and
$R_-=\theta_-+m\Delta_A$. Substitution produces the last two terms of
\eqref{ms:simultaneous-gap-formula} and the deterministic expression
\[
 m\{(\sigma-1)\bar B^2+(\sigma+1)\bar A^2-6\bar A\bar B\}
 +6m\{\bar A\bar B-AB+B\Delta_B+A\Delta_A\}.
\]
The second brace equals
$(A+B)(\Delta_A+\Delta_B)+\Delta_A\Delta_B$.
For the first, use $m^2=(1-\theta)pq$ and the exact identity
\[
 (\sigma+1)p^2+(\sigma-1)q^2-6pq=(6+2\sigma)(p-p_*)^2.
\]
This proves the displayed inequality, with one residual paying for all
the nonnegative costs simultaneously.
\end{proof}

In count coordinates the inward displacements are
\[
 a_0=p-A/\ell=\Delta_A/\ell,\quad
 a_1=q-B/\ell=\Delta_B/\ell,
 \quad d_Y=a_0+a_1=(\Delta_A+\Delta_B)/\ell.
\]
Thus $U_0\le a_0$ and $U_1\ge-a_1$. The physical displacement
$\Delta_A+\Delta_B$ and the count displacement $d_Y$ differ by $\ell$.

\begin{theorem}[Weighted joint resource]
\label{ms:joint-resource}
Suppose $n\ge N>0$ and established bounds for this same law give
\[
 nD\le Nd_D,\quad \theta\le\Theta<1,\quad pq\le v_*,\quad
 D\ge b+6mS_*(\Delta_A+\Delta_B)+K\theta,
\]
where $b\ge0$, $S_*>0$, $K>0$ and $v_*>0$. Put $r=d_D-b$.
If $r<0$ there is no such law. If $r=0$, then $d_Y=\lambda=0$.
If $r>0$, define
\[
 d_* =\frac{r}{6(1-\Theta)S_*},\qquad
 \Lambda=\frac{Nv_*r}{K(1-\Theta)}.
\]
Then
\begin{equation}
 \frac{d_Y}{d_*}+\frac\lambda\Lambda\le1.
\label{ms:joint-plane}
\end{equation}
\end{theorem}
\begin{proof}
Since $m\ell=1-\theta$ and
$n\theta=(1-\theta)\lambda/(pq)$,
\begin{align*}
 Nd_D&\ge nD\ge nb+6n(1-\theta)S_*d_Y
                          +K(1-\theta)\lambda/(pq)\\
 &\ge Nb+6N(1-\Theta)S_*d_Y+K(1-\Theta)\lambda/v_*.
\end{align*}
All weakened coefficients multiply nonnegative quantities. Subtracting
$Nb$ proves the impossible and zero cases without division. For $r>0$,
division by $Nr$ proves \eqref{ms:joint-plane}.
\end{proof}

For instance, on a closed probability interval $I$, established endpoint
floors $A\ge A_l>0$, $B\ge B_l>0$ and $\theta\le\Theta$ supply
\[
 S_*=A_l+B_l,\qquad K=\min\{3(\sigma+1)B_l,3(\sigma+3)A_l\},
\]
\[
 b=\frac{(6+2\sigma)(1-\Theta)^{3/2}}{\sqrt{v_*}}
          \operatorname{dist}(p_*,I)^2,
 \qquad v_*\ge\max_{p\in I}p(1-p).
\]
These choices follow from Lemma~\ref{ms:simultaneous-gap} by dropping the
nonnegative product $6m\Delta_A\Delta_B$. Rational lower enclosures for
$b,K,S_*$ and an upper enclosure for $v_*$ preserve the implication.
The preceding theorem assumes the weighted statement $nD\le Nd_D$;
an unweighted $D\le d_D$ cannot be multiplied by $N$ when $n>N$.

On a band $L\le\lambda\le U\le\Lambda$, the joint plane gives
$d_Y\le d_*(1-L/\Lambda)$. Either positive intercept may be enlarged
outward. A separate smaller marginal noise bound does not, by itself,
replace $\Lambda$ in this joint plane. The two intercepts must still come
from a proved simultaneous inequality.
